\documentclass[10pt,twocolumn,a4paper]{article}

% Pacotes
\usepackage[utf8]{inputenc}
\usepackage[T1]{fontenc}
\usepackage[brazil]{babel}
\usepackage{geometry}
\usepackage{graphicx}
\usepackage{amsmath, amssymb}
\usepackage{booktabs}
\usepackage{cite}
\usepackage{hyperref}
\usepackage{listings}
\usepackage{xcolor}

% Configuração do listings para Java
\lstdefinestyle{javastyle}{
    language=Java,
    basicstyle=\ttfamily\footnotesize,
    keywordstyle=\color{blue}\bfseries,
    commentstyle=\color{gray}\itshape,
    stringstyle=\color{orange},
    numberstyle=\tiny\color{gray},
    numbers=left,
    stepnumber=1,
    numbersep=5pt,
    backgroundcolor=\color{gray!8},
    frame=single,
    rulecolor=\color{gray!40},
    breaklines=true,
    breakatwhitespace=false,
    captionpos=b,
    tabsize=4,
    showstringspaces=false,
    xleftmargin=8pt,
    xrightmargin=4pt,
}
\lstset{style=javastyle}

% Margens
\geometry{left=2cm, right=2cm, top=2cm, bottom=2cm}
\setlength{\columnsep}{0.7cm}

% Título
\title{\textbf{Implementação de um chat por protocolo TCP}\\
\large Sistemas Distribuídos / Programação Paralela -- UTFPR}

\author{
    Edilson Daniel de Souza \\
    \texttt{edilsonsouza@alunos.utfpr.edu.br}
    \and
    Samuel Minuzzo Pereira \\
    \texttt{samuelpereira@alunos.utfpr.edu.br}
}

\date{\today}

\begin{document}

\maketitle

% Resumo
\begin{abstract}
Este trabalho apresenta a implementação de um sistema de chat em tempo real
utilizando o protocolo de comunicação TCP em Java. A aplicação é composta por
um servidor centralizado e múltiplos clientes com interface gráfica, permitindo
troca de mensagens públicas e privadas entre usuários conectados
simultaneamente. A arquitetura adota \textit{threads} para o tratamento
concorrente de conexões no servidor e para a recepção assíncrona de mensagens
no cliente, garantindo responsividade e estabilidade da aplicação.
\end{abstract}

\textbf{Palavras-chave:} Programação Paralela, Threads, TCP, Comunicação.

% ------------------------------------------------

\section{Introdução}
A comunicação em rede é um dos pilares da computação moderna. Protocolos como
o TCP (\textit{Transmission Control Protocol}) garantem entrega confiável e
ordenada de dados entre processos distribuídos, sendo amplamente utilizados em
aplicações onde a entrega da mensagem é indispensável.

Este trabalho tem como motivação a aplicação prática dos conceitos de
programação paralela e distribuída, explorando o uso de \textit{sockets} TCP e
\textit{threads} para construir um sistema funcional de chat multiusuário em
Java.

O objetivo principal é desenvolver uma aplicação cliente-servidor que permita:
\begin{itemize}
    \item Comunicação simultânea entre múltiplos usuários;
    \item Troca de mensagens públicas (\textit{broadcast}) e privadas;
    \item Interface gráfica intuitiva com lista de usuários online;
    \item Identificação dos usuários e registro de data/hora nas mensagens.
\end{itemize}

% ------------------------------------------------

\section{Descrição do Problema}
O problema consiste em implementar um sistema de chat baseado no modelo
cliente-servidor utilizando \textit{sockets} TCP. O servidor deve aceitar
conexões de múltiplos clientes de forma concorrente, receber mensagens e
encaminhá-las ao(s) destinatário(s) correto(s).

\textbf{Entrada:} mensagens de texto digitadas pelo usuário, podendo ser
direcionadas a todos os participantes ou a um usuário específico.

\textbf{Saída:} exibição das mensagens recebidas na interface gráfica de cada
cliente, formatadas com remetente, destinatário e carimbo de data/hora.

\textbf{Restrições:}
\begin{itemize}
    \item Uso exclusivo de sockets TCP para comunicação;
    \item Objetos serializáveis para tráfego de mensagens;
    \item Tratamento adequado de desconexões e exceções.
\end{itemize}

% ------------------------------------------------

\section{Implementações}
O sistema é composto por três classes principais: \texttt{Server},
\texttt{Client} e \texttt{Mensagem}, todas pertencentes ao pacote \texttt{chat}.

% ---
\subsection{Classe \texttt{Mensagem}}

A classe \texttt{Mensagem} implementa \texttt{Serializable}, permitindo que
objetos sejam transmitidos diretamente pelos \textit{streams} de objetos do
Java. Ela encapsula os campos \texttt{remetente}, \texttt{destinatario},
\texttt{conteudo} e \texttt{horario} (do tipo \texttt{LocalDateTime}).

A convenção adotada é: quando \texttt{destinatario == null}, a mensagem é
tratada como \textit{broadcast}; caso contrário, é uma mensagem privada.

O método \texttt{toString()} formata a saída no padrão exigido:

\begin{lstlisting}[caption={toString() da classe Mensagem}]
@Override
public String toString() {
    String destino = (destinatario == null)
        ? "Todos" : destinatario;
    return "[" + horario.format(FORMATTER) + "] "
        + remetente + " -> " + destino
        + ": " + conteudo;
}
\end{lstlisting}

O formatador utiliza o padrão \texttt{dd/MM/yyyy HH:mm:ss}, garantindo que
cada mensagem exibida traga o instante exato de criação.

% ---
\subsection{Classe \texttt{Server}}

O servidor é iniciado com um \texttt{ServerSocket} na porta 1234 e entra em
laço infinito aguardando conexões via \texttt{accept()}. Para cada cliente
aceito, é criado um \texttt{ClientHandler} — que implementa \texttt{Runnable}
— iniciado em uma \textit{thread} separada.

A lista de clientes conectados é mantida em uma
\texttt{CopyOnWriteArrayList<ClientHandler>}, garantindo segurança em ambientes
concorrentes sem necessidade de sincronização explícita.

\subsubsection{Construtor do \texttt{ClientHandler}}

O construtor do \texttt{ClientHandler} é responsável por todo o
\textit{handshake} inicial: cria os \textit{streams}, lê o nome enviado pelo
cliente e rejeita a conexão caso o nome seja vazio ou já esteja em uso.

\begin{lstlisting}[caption={Handshake no ClientHandler}]
public ClientHandler(Socket socket)
        throws IOException, ClassNotFoundException {
    this.socket = socket;
    out = new ObjectOutputStream(
        socket.getOutputStream());
    out.flush(); // envia cabecalho antes de ler
    in  = new ObjectInputStream(
        socket.getInputStream());
    nome = ((String) in.readObject()).trim();
    if (nome.isEmpty())
        throw new IOException("Nome invalido");
    if (Server.nomeJaExiste(nome)) {
        out.writeObject(new Mensagem("SERVIDOR",
            nome, "Nome ja esta em uso."));
        out.flush();
        socket.close();
        throw new IOException("Nome duplicado");
    }
}
\end{lstlisting}

O \texttt{ObjectOutputStream} é criado e descarregado antes do
\texttt{ObjectInputStream} para evitar \textit{deadlock} durante o
\textit{handshake} (ver Seção~\ref{sec:dificuldades}).

\subsubsection{Broadcast e mensagem privada}

\begin{lstlisting}[caption={Broadcast e envio privado}]
public void broadcast(Mensagem msg) {
    for (ClientHandler c : clientes)
        enviar(c, msg);
}

private void enviarPrivado(Mensagem msg) {
    boolean encontrado = false;
    for (ClientHandler c : clientes) {
        if (c.nome.equals(msg.getDestinatario())
         || c.nome.equals(msg.getRemetente())) {
            enviar(c, msg);
            encontrado = true;
        }
    }
    if (!encontrado)
        enviar(this, new Mensagem("SERVIDOR", nome,
            "Usuario '" + msg.getDestinatario()
            + "' nao encontrado."));
}
\end{lstlisting}

O método \texttt{enviarPrivado} entrega a mensagem tanto ao destinatário quanto
ao próprio remetente — garantindo que este veja a mensagem enviada na própria
janela — e notifica o remetente caso o destinatário não esteja conectado.

\subsubsection{Atualização da lista online}

Sempre que um cliente entra ou sai do chat, o método
\texttt{atualizarUsuarios()} transmite a todos uma \texttt{String} no formato
\texttt{ONLINE:usuario1,usuario2,...}, que o lado cliente usa para atualizar o
painel lateral.

\begin{lstlisting}[caption={Envio da lista de usuários online}]
public void atualizarUsuarios() {
    StringBuilder sb = new StringBuilder();
    for (ClientHandler c : clientes)
        sb.append(c.nome).append(",");
    String lista = "ONLINE:" + sb;
    for (ClientHandler c : clientes)
        enviar(c, lista);
}
\end{lstlisting}

% ---
\subsection{Classe \texttt{Client} e Interface Gráfica}

O cliente possui interface gráfica construída com Swing. Ao iniciar, uma caixa
de diálogo solicita o nome do usuário e permite configurar o endereço do
servidor — localhost ou IP personalizado —, eliminando a necessidade de editar
o código-fonte para conexões em rede local.

\subsubsection{Estrutura da janela principal}

A janela principal é organizada em três regiões com \texttt{BorderLayout}:

\begin{itemize}
    \item \textbf{Painel esquerdo (\texttt{WEST}):} \texttt{JList} de usuários
          online com seleção por clique e botão ``Todos'' para retornar ao
          \textit{broadcast};
    \item \textbf{Centro (\texttt{CENTER}):} \texttt{JTextArea} somente leitura
          envolto em \texttt{JScrollPane} para exibição do histórico de
          mensagens;
    \item \textbf{Rodapé (\texttt{SOUTH}):} \texttt{JLabel} com o destinatário
          atual, \texttt{JTextField} de digitação e \texttt{JButton} de envio.
\end{itemize}

\begin{lstlisting}[caption={Montagem do painel inferior (rodapé)}]
JPanel painelInferior =
    new JPanel(new BorderLayout());
labelDestinatario = new JLabel("  Para: Todos  ");
campoMensagem     = new JTextField();
botaoEnviar       = new JButton("Enviar");
painelInferior.add(
    labelDestinatario, BorderLayout.WEST);
painelInferior.add(
    campoMensagem,     BorderLayout.CENTER);
painelInferior.add(
    botaoEnviar,       BorderLayout.EAST);
add(painelInferior, BorderLayout.SOUTH);
\end{lstlisting}

\subsubsection{Seleção de destinatário via clique}

O clique em um nome da lista lateral altera o destinatário ativo e atualiza
o rótulo no rodapé. O botão ``Todos'' volta o destinatário para
\texttt{null}, reativando o \textit{broadcast}.

\begin{lstlisting}[caption={Listener da lista de usuários}]
listaUsuarios.addMouseListener(
  new MouseAdapter() {
    @Override
    public void mouseClicked(MouseEvent e) {
        String sel =
            listaUsuarios.getSelectedValue();
        if (sel != null) {
            destinatario = sel;
            labelDestinatario.setText(
                "  Para: " + destinatario + "  ");
        }
    }
});
\end{lstlisting}

\subsubsection{Thread receptora (\texttt{ReceiverThread})}

A recepção de mensagens é delegada à classe interna \texttt{ReceiverThread},
que lê objetos do \texttt{ObjectInputStream} em laço contínuo. Quando recebe
uma \texttt{String} com prefixo \texttt{ONLINE:}, delega a atualização da lista
para o EDT via \texttt{SwingUtilities.invokeLater}. Quando recebe um objeto
\texttt{Mensagem}, adiciona o texto à área de chat.

\begin{lstlisting}[caption={Loop principal do ReceiverThread}]
while (true) {
    Object obj = in.readObject();
    if (obj instanceof String) {
        String str = (String) obj;
        if (str.startsWith("ONLINE:")) {
            synchronized (ReceiverThread.class) {
                listaUsuarios = str.substring(7);
            }
            SwingUtilities.invokeLater(() ->
                client.atualizarListaUsuarios(
                    getListaUsuarios()));
        }
    } else if (obj instanceof Mensagem) {
        Mensagem msg = (Mensagem) obj;
        if (msg.getRemetente().equals("SERVIDOR")
         && msg.getConteudo()
               .contains("Nome ja esta em uso")) {
            client.reconectarComNovoNome();
            break;
        }
        SwingUtilities.invokeLater(() ->
            areaChat.append(msg + "\n"));
    }
}
\end{lstlisting}

\subsection{Comandos Suportados}

\begin{itemize}
    \item \texttt{/usuarios} — solicita ao servidor a lista de clientes
          conectados;
    \item \texttt{/privado :<usuário> :<mensagem>} — envia mensagem privada
          via linha de comando;
    \item Clique na lista lateral — seleciona destinatário privado diretamente
          pela interface gráfica.
\end{itemize}

% ------------------------------------------------

\section{Aspectos de Concorrência}

A concorrência é tratada em dois níveis:

\textbf{No servidor,} cada cliente é atendido por uma \textit{thread}
independente (\texttt{ClientHandler}). A estrutura
\texttt{CopyOnWriteArrayList} permite iteração segura sobre a lista de clientes
durante operações de \textit{broadcast}, sem bloqueios explícitos.

\textbf{No cliente,} a \texttt{ReceiverThread} opera independentemente da
\textit{thread} da interface gráfica (EDT). O acesso à variável estática
\texttt{listaUsuarios} dentro de \texttt{ReceiverThread} é protegido por bloco
\texttt{synchronized}, evitando condições de corrida.

A ordem de inicialização dos \textit{streams} — \texttt{ObjectOutputStream}
antes de \texttt{ObjectInputStream} — é respeitada em ambos os lados da
conexão, prevenindo \textit{deadlock} durante o \textit{handshake} inicial.

% ------------------------------------------------

\section{Tratamento de Erros}

A aplicação trata os seguintes cenários de falha:

\begin{itemize}
    \item \textbf{Desconexão de cliente:} a exceção de I/O capturada no laço do
          \texttt{ClientHandler} aciona \texttt{fecharConexao()}, que remove o
          cliente da lista, notifica os demais e atualiza a lista online;
    \item \textbf{Destinatário inexistente:} o servidor notifica o remetente
          com mensagem de aviso;
    \item \textbf{Comando mal formatado:} o cliente exibe mensagem de uso
          correto na área de chat sem encerrar a aplicação;
    \item \textbf{Falha de conexão no cliente:} exibição de diálogo de erro
          via \texttt{JOptionPane}.
\end{itemize}

Todos os blocos \texttt{catch} utilizam \textit{multi-catch} onde aplicável
(\texttt{IOException | ClassNotFoundException}), e blocos \texttt{finally}
garantem o registro de mensagens de diagnóstico no console.

% ------------------------------------------------

\section{Instruções de Execução}

\subsection{Compilação}

O projeto deve ser compilado a partir do diretório pai da pasta \texttt{chat/}:

\begin{verbatim}
javac chat/*.java
\end{verbatim}

\subsection{Execução}

\begin{verbatim}
# Iniciar o servidor
java chat.Server

# Iniciar um cliente (em outro terminal)
java chat.Client
\end{verbatim}

Ao iniciar o cliente, uma janela de configuração permite definir o nome de
usuário e o endereço do servidor. Múltiplas instâncias do cliente podem ser
abertas simultaneamente.

% ------------------------------------------------

\section{Dificuldades Encontradas}
\label{sec:dificuldades}

\begin{itemize}
    \item \textbf{Deadlock na inicialização dos streams:} o Java bloqueia na
          criação do \texttt{ObjectInputStream} até receber o cabeçalho do
          \texttt{ObjectOutputStream} do lado oposto. A solução foi garantir
          que o \texttt{ObjectOutputStream} seja sempre criado e descarregado
          (\texttt{flush()}) antes do \texttt{ObjectInputStream} em ambos os
          lados.

    \item \textbf{Atualização da interface a partir de threads:} chamadas
          diretas a componentes Swing fora do EDT causam comportamentos
          imprevisíveis. O uso de \texttt{SwingUtilities.invokeLater} resolveu
          o problema de forma consistente.

    \item \textbf{Sincronização da lista de usuários:} o cache estático em
          \texttt{ReceiverThread} precisou de bloco \texttt{synchronized} para
          evitar leituras inconsistentes entre a thread receptora e o EDT.
\end{itemize}


\section{Resultados}

\subsection{Comportamento da Interface Gráfica}

A interface gráfica mostrou-se responsiva em todos os testes. A lista de
usuários online foi atualizada automaticamente a cada entrada ou saída de
cliente, sem atrasos perceptíveis. O rótulo de destinatário no rodapé manteve
o estado correto mesmo após a saída de um usuário selecionado — o sistema
voltou automaticamente ao modo \textit{broadcast}.

\subsection{Consistência das Mensagens}

As mensagens foram entregues na ordem correta em todos os testes. O carimbo
de data/hora refletiu com precisão o instante de criação de cada objeto
\texttt{Mensagem} no cliente remetente. Em mensagens privadas, a entrega foi
confirmada tanto no remetente quanto no destinatário, sem vazamento para outros
clientes conectados.

% ------------------------------------------------

\section{Discussão}

\subsection{Decisões de Arquitetura}

A escolha da \texttt{CopyOnWriteArrayList} para armazenar os
\texttt{ClientHandler} no servidor simplificou o código ao eliminar blocos
\texttt{synchronized} explícitos nas operações de \textit{broadcast}.
O custo de cópia interna da estrutura a cada escrita é aceitável dado que
adições e remoções ocorrem raramente (apenas em conexões e desconexões),
enquanto iterações de leitura são frequentes a cada mensagem enviada.

\subsection{Escalabilidade}

O modelo de uma \textit{thread} por cliente é simples e eficaz para cenários
com dezenas de usuários simultâneos. Para escalas maiores (centenas ou
milhares de conexões), a abordagem com \texttt{NIO} (Non-blocking I/O) e um
pool de \textit{threads} seria mais adequada, pois reduziria o overhead de
criação e chaveamento de \textit{threads}.

\subsection{Segurança e Robustez}

A aplicação não implementa criptografia: as mensagens trafegam em texto claro
pelo socket TCP. Para uso em redes não confiáveis, seria necessário adotar
\texttt{SSLSocket} em substituição ao \texttt{Socket} padrão. A validação de
nomes se limita à verificação de duplicidade e nome vazio, não restringindo o uso
caracteres especiais.

\subsection{Experiência de Uso}

O diálogo de configuração inicial elimina a necessidade de editar código-fonte
para alterar o IP do servidor, tornando a aplicação portável para uso em
laboratório. A seleção de destinatário por clique é mais fluida do que o
uso do comando textual \texttt{/privado}, demonstrando a vantagem da abordagem
híbrida (GUI + comandos) adotada.

% ------------------------------------------------

\section{Conclusão}

Este trabalho demonstrou a viabilidade de construir um sistema de chat
multiusuário em tempo real utilizando exclusivamente recursos nativos do Java:
\textit{sockets} TCP, serialização de objetos, \textit{threads} e a biblioteca
Swing para a interface gráfica.

Os principais objetivos foram alcançados: múltiplos clientes se comunicam
simultaneamente, mensagens públicas e privadas são entregues corretamente, a
lista de usuários online é mantida atualizada de forma automática e as
desconexões são tratadas sem comprometer os demais participantes.

Os aspectos de concorrência — uso de \texttt{CopyOnWriteArrayList} no servidor
e separação entre \texttt{ReceiverThread} e o EDT no cliente — provaram-se
suficientes para garantir estabilidade e ausência de condições de corrida nos
testes realizados.

Como trabalhos futuros, destacam-se: adoção de \texttt{SSLSocket} para
criptografia do tráfego, migração para NIO com pool de \textit{threads} para
suporte a maior escala, persistência do histórico de mensagens em banco de
dados, e suporte a salas de chat (grupos) com controle de permissões.

% ------------------------------------------------

\begin{thebibliography}{99}

\bibitem{tanenbaum}
TANENBAUM, A. S.; VAN STEEN, M.
\textit{Distributed Systems: Principles and Paradigms}. 3. ed.
Pearson, 2017.

\bibitem{oracle-sockets}
Oracle Corporation.
\textit{All About Sockets} -- The Java\texttrademark\ Tutorials.
Disponível em: \url{https://docs.oracle.com/javase/tutorial/networking/sockets/}.
Acesso em: 2025.

\bibitem{oracle-swing}
Oracle Corporation.
\textit{Creating a GUI With Swing} -- The Java\texttrademark\ Tutorials.
Disponível em: \url{https://docs.oracle.com/javase/tutorial/uiswing/}.
Acesso em: 2025.

\bibitem{goetz}
GOETZ, B. et al.
\textit{Java Concurrency in Practice}. Addison-Wesley, 2006.

\end{thebibliography}

\end{document}